\section{Ergodicity and the invariant $\sigma$-algebra}
\label{sec:ergodic}

This section establishes the ergodic-theoretic input to the proof
of Theorem~\ref{thm:empirical_invariance}: the shift~$T$ on the
path space $\Omega$ is ergodic with respect to the empirical
measure $\pi_m^{\mathrm{emp}}$, the invariant $\sigma$-algebra is
trivial, and the geometry $G(\mu)$ is therefore almost surely
constant. The proofs use only the spectral gap of $T_m$ and the
Birkhoff--Hopf ergodic theorem; all the technical heavy lifting
takes place at the level of the gap field $\{g_n\}$ and is
inherited from chapters~7 and~8 of the
monograph~\cite{PT_Monograph}.

\subsection{The shift dynamical system on trajectories}
\label{ssec:shift}

Recall that the path space $\Omega = \mathbb{N}^{\mathbb{N}}$ is
equipped with the cylindrical $\sigma$-algebra $\mathcal{F}$ and
the canonical shift $T : \Omega \to \Omega$,
$T(\omega)_n = \omega_{n+1}$. The canonical trajectory
$\omega^\ast = (g_1, g_2, g_3, \ldots)$ identifies $\Omega$ with
the orbit closure of $\omega^\ast$ under $T$. For each squarefree
integer $m$ and each $n \in \mathbb{N}$, the modular coordinate
$r_n^{(m)} = g_n \bmod m$ defines a finite-valued random variable
on $\Omega$; the joint process $(r_n^{(m)})_{n \geq 1}$ is a
Markov chain with transition matrix $T_m$, by the definition of
the sieve.

For each prime~$p$, the trajectory $\sigma$-algebra
$\mathcal{A}_p = \sigma(\{r_n^{(p)} : n \in \mathbb{N}\})$
introduced in Eq.~\eqref{eq:Ap_def} is the smallest
sub-$\sigma$-algebra of $\mathcal{F}$ that contains every
finite-dimensional cylinder built from the mod-$p$ residue
coordinates. By construction, $\mathcal{A}_p$ is $T$-invariant in
the weak sense $T^{-1} \mathcal{A}_p \subseteq \mathcal{A}_p$,
but not in the strict sense $T^{-1} \mathcal{A}_p =
\mathcal{A}_p$; we will not need this distinction. What we will
need is the shift-invariant $\sigma$-algebra
\begin{equation}
  \mathcal{I}
  \;:=\;
  \bigl\{ A \in \mathcal{F} : T^{-1} A = A \bigr\}
  \;\subset\; \mathcal{F},
  \label{eq:I_def}
\end{equation}
the family of measurable subsets of $\Omega$ that are exactly
invariant under the shift.

\subsection{Ergodicity of the prime-sieve flow}
\label{ssec:ergodicity}

\begin{proposition}[Sieve Ergodicity]
\label{prop:ergodic}
The shift $T$ on $\Omega$ is ergodic with respect to
$\pi_m^{\mathrm{emp}}$.
\end{proposition}

\begin{proof}[Proof sketch]
The argument has two ingredients. The first is the spectral-gap
estimate of Theorem~T4 in the
monograph~\cite[Ch.~7,~Thm.~T4]{PT_Monograph}: the second-largest
eigenvalue of the per-prime transfer matrix $T_p$ on the active
subspace satisfies $|\lambda^{(p)}_1| = 1/(p-2) < 1$ for $p \geq 5$
and $|\lambda^{(3)}_1| = 1$ with antiperiodic structure. The
tensor decomposition~\eqref{eq:Tm_tensor} propagates the spectral
gap to the joint matrix $T_m$, so that on the active subspace
$\mathcal{S}_m$ the difference
$\|T_m^n \nu - \pi_m^{\mathrm{Ruelle}}\|_{\mathrm{TV}}$ decays
geometrically for any initial distribution~$\nu$ with positive
support on $\mathcal{S}_m$.

The second ingredient is the Mertens-type estimate of
Section~7.4 of the monograph~\cite{PT_Monograph}, which
establishes that the empirical measure $\pi_m^{\mathrm{emp}}$ of
the canonical trajectory $\omega^\ast$ converges to
$\pi_m^{\mathrm{Ruelle}}$ in total variation:
$\|\pi_m^{\mathrm{emp}} - \pi_m^{\mathrm{Ruelle}}\|_{\mathrm{TV}} \to 0$
as the trajectory is extended to infinity, up to a residual term
controlled by the multiplicative constant of Mertens' third
theorem~\cite[Ch.~22]{Davenport2000}.

These two ingredients together imply that
$(\Omega, \mathcal{F}, \pi_m^{\mathrm{emp}}, T)$ is a measure-preserving
dynamical system with the same invariant statistics as the
irreducible aperiodic Markov chain on $\mathcal{S}_m$ defined by
$T_m$. The Birkhoff--Hopf ergodic
theorem~\cite[Thm.~1.6]{Walters1982} then applies, and gives
ergodicity of $T$ with respect to $\pi_m^{\mathrm{emp}}$.
\end{proof}

\begin{remark}[On the precise sense of ergodicity]
\label{rmk:precise_ergodicity}
Proposition~\ref{prop:ergodic} requires a slight reading of the
notion of ergodicity, because $\pi_m^{\mathrm{emp}}$ is defined
through a long-time average rather than an a priori invariant
measure. The standard fix~\cite[Ch.~1.6,~Remark~1]{Walters1982} is
to verify that $\pi_m^{\mathrm{emp}}$ is shift-invariant
(immediate from the definition~\eqref{eq:pi_emp_def}) and that
every shift-invariant subset of $\Omega$ has measure $0$ or $1$
under $\pi_m^{\mathrm{emp}}$. The latter property follows from the
spectral-gap estimate above, by a standard Birkhoff--Hopf-style
argument that we do not reproduce here; the canonical reference is
chapter~6 of Krengel's monograph~\cite{KrengelBook}.
\end{remark}

\subsection{Triviality of the invariant $\sigma$-algebra}
\label{ssec:trivial_inv}

\begin{corollary}[Invariant Triviality (Birkhoff--Hopf)]
\label{cor:trivial}
The $T$-invariant $\sigma$-algebra
$\mathcal{I} = \{A : T^{-1}A = A\}$ is trivial modulo
$\pi_m^{\mathrm{emp}}$:
$\mathcal{I} = \{\emptyset, \Omega\} \pmod{\pi_m^{\mathrm{emp}}}$.
\end{corollary}

\begin{proof}
This is the standard equivalence between ergodicity and triviality
of the invariant $\sigma$-algebra; see Theorem~1.6 (iv)
of~\cite{Walters1982}. Given Proposition~\ref{prop:ergodic}, the
result follows immediately.
\end{proof}

\subsection{Spectral measurability $\Rightarrow$ a.s.\ constancy
of $G(\mu)$}
\label{ssec:G_as_constant}

\begin{corollary}[Geometric A.S.\ Constancy]
\label{cor:G_constant}
For every $\mu \in [\mu_c, \infty)$, the geometry $G(\mu)$ is
$\pi_m^{\mathrm{emp}}$-almost-surely constant.
\end{corollary}

\begin{proof}
By Lemma~\ref{lem:G_measurable}, $G(\mu) \in
\sigma(\mathrm{Spec}(T_m)) \subseteq \mathcal{I}$. By
Corollary~\ref{cor:trivial}, any $\mathcal{I}$-measurable function
is $\pi_m^{\mathrm{emp}}$-almost-surely constant.
\end{proof}

\paragraph{Why this is the substantive content.}
At a first reading, the conclusion of Corollary~\ref{cor:G_constant}
might appear vacuous: $G(\mu)$ is, after all, a deterministic
function of $\mu$, and ``almost surely constant'' looks like an
empty statement about deterministic objects. The substance lies in
the role that $G(\mu)$ plays in the conditional mutual
information $I(\Phi_p; \Phi_q \mid G(\mu))$: even though
$G(\mu)$ is deterministic as a function of $\mu$, the random
variable $G$ in the conditioning is being interpreted as a
$\mathcal{F}$-measurable random variable on the path space $\Omega$,
and its $\Omega$-measurable image is what matters for the
information-theoretic identity used in Step~4 of the proof of
Theorem~\ref{thm:empirical_invariance}. The triviality of
$\mathcal{I}$ modulo $\pi_m^{\mathrm{emp}}$ collapses this
$\Omega$-measurable image to a single point, after which the
identity
$I(\Phi_p; \Phi_q \mid G) = I(\Phi_p; \Phi_q)$ becomes available.
The point of the corollary is not the deterministic statement that
$G(\mu)$ is a fixed matrix; it is the measure-theoretic statement
that, viewed as a random variable on the path space, $G(\mu)$
contributes nothing to the joint information content of the modular
fibres.
